\clearpage

% \section{TODO}

% \begin{enumerate}
% \item RIPr and convex optimization
%     \item Formal on variable block distribution, filtrations, and sigma-algebra, $n_a, n_b$
%     \item Add manual data for verification
%     \item Benchmark \E-values calculated from UI, Turners, calibrator (from Fisher), conditional (1-shot) \\
%     I called the conditional \E-variables looking at the entire sequence as one-shot, maybe a better in Hao's
%     \item Draw a picture of $\gamma = 0.18$ is not the best, 
%     this could be done in old code, right?
%     Also, $0.18$ was calculated under $\theta_a = \theta_b$, prove that 
%     both numerically and analytically
%     \item Pseudo-algorithm for $S_{\textsc{cond}}$, 2 cond \E: one takes
%     a single data block at a time 
%     while the other take the whole sequence of $m$ points.
%     The former is in the k-stream paper and later is in the general case paper.
%     The later is optimal: same \E-power as RIPr asymptotically 
%     but former is a bit worse. (After discussion)
%     \item Draw a picture of $S_{\textsc{mix}}$ vs $S_{\textsc{cond}}$
%     \item Write down growth rate difference in different cases (Draw also)
%     \item Seq-RIPr, Bayesian Predictive Model, Not in exp. family
%     \item Revisit the degenerate of the Prior in Theorem 2 \cite{turnerExactAnytimevalidConfidence2023}
%     \item Carathéodory’s Theorem
%     \item \sout{Write down $\hypN$ and $\hypA$ explicitly, especially for the boundary conditions}
%     \item \sout{How to calculate RIPr: first derivative being zero or direct minimisation of KL?}
%     \item \sout{Pseudo-algorithm for the calculation of CS}
%     \item \sout{Why bootstrapping? to get the SE for stopping times}
%     \item \sout{Pseudo-algorithm of calculate when $\E \geq 1/\alpha$}
% \end{enumerate}

% \subsection{Problem with previous methods}
% \begin{enumerate}
%     \item \verb|pilot = TRUE| does nothing
%     \item Possible wrong calculation with highly unbalanced group design $n_a \ll n_b$
% \end{enumerate}

% \subsection{Must have}
% \begin{enumerate}
%     \item a function to calculate \E\ without simulation:\\
%         balanced and unbalanced, just \verb|calculateSequential2x2E| maybe?
%     \item calculate $\CS{m}$ for each $m$
%     \item design: calculate \verb|nPlan| with $\delta$ and $1 - \mu$, same as calculate power $1 - \mu$ with \verb|nPlan| and $\delta$
%     \item design: calculate $\delta$ with $1 - \mu$ and \verb|nPlan|
% \end{enumerate}

% \subsection{Nice to have}

% \begin{enumerate}
%     \item 2x2 in python in \url{https://github.com/jakorostami/expectation}
%     \item calculate conditional \E-variable
%     \item fast way to compute RIPr in each specified settings, especially for CS
%     \item update GitHub CI/CD
%     \item pre-commit
%     \item unify \texttt{R} style
% \end{enumerate}

% \clearpage


\section{SAVI Contingency Table}


\subsection{Setup and notation}

We denote the Kullback–Leibler divergence by $\KL{P}{Q} = \int \log(\dv*{P}{Q}) \dd{P}$
where $\dv*{P}{Q}$ is the Radon-Nikodym derivative of $P$ with respect to $Q$ and
the integral is taken on the same probability space (not shown here).
We will assume that the measures discussed here all have full support.

We start with a binary random variable $U$
following Bernoulli distribution with parameter $\mu$,
here $\mu$ is also the mean in mean-value parameterization.
The density reads
\begin{align*}
	p_{\mu}(U)
	 & := \mu^U \cdot (1 - \mu)^{1-U}, \quad U \in \{0,1\}, \quad \mu \in [0,1]
\end{align*}

By the standard properties of regular exponential families,
we parameterize $\Pc$ in terms of the mean-value parameter space $\Mt_p$ and write
\begin{align*}
	\Pc = \{P_\mu: \mu \in \Mt_p\}, \quad \Mt_p = [0,1].
\end{align*}

The canonical parameterization for any $\mu \in \Mt_p$ is defined as
\begin{align*}
	p_{\beta}(U) = \frac{1}{Z_p(\beta)} \cdot \exp(\beta \cdot U) \cdot p_{\mu}(U),
\end{align*}
where $Z_p(\beta)$ is the normalizing constance and the canonical parameter space is
defined as $\Bt_p = \{\beta: Z_p(\beta) < \infty \} = \Rf$.
% \beta = \log(\frac{\mu}{1-\mu})

Moving on to the 2-stream examples with different group size $n_a, n_b$.
At each time point $i$, $n_a$ and $n_b$ copies of Bernoulli variables are collected,
following the same distribution:
\begin{align*}
	U_a & = (U_{a,1}, U_{a,2},\cdots,U_{a,n_a})^T \sim \mathrm{Ber}(\bm{\mu}_a) \\
	U_b & = (U_{b,1}, U_{b,2},\cdots,U_{b,n_n})^T \sim \mathrm{Ber}(\bm{\mu}_b)
\end{align*}
The sufficient statistics vector is denoted by $X = (\sum_j^{n_a}U_{a,j}, \sum_j^{n_b}U_{b,j})$.
$\Pc$ is extended by the i.i.d. assumption to
a set of distributions for random process $U_{(1)}, U_{(2)}, \dots$
for $i \geq 1$, i.i.d. copies of $U_{(i)}=(U_{(i),a}, U_{(i),b})$.
Using superscript, we denote the random process up to time $i$ by
$U^{(n)} := (U_{(1)}, U_{(2)}, \dots U_{(n)})$
and $X^{(n)} := (X_{(1)}, X_{(2)}, \dots X_{(n)})$ for the sufficient statistics respectively.

The \emph{alternative} we consider first $\Qc = \{Q\}$ is simple with some oracle
$\bm{\mu}^* = (\bm{\mu}_a^*, \bm{\mu}_b^*)$.
\begin{align*}
	\hypA: U_a \sim \mathrm{Ber}(\bm{\mu}_a^*), U_b \sim \mathrm{Ber}(\bm{\mu}_b^*) \quad
	\text{independent}.
\end{align*}

Under mean-value parameterization:
\begin{align*}
	\hypA: X_a \sim P_{\bm{\mu}_a^*}, X_b \sim P_{\bm{\mu}_b^*} \quad
	\text{independent}.
\end{align*}

The null distribution $\Pc$ is the set of distributions where group $a$ and $b$
are equipped with the same parameter and hence composite:
\begin{align*}
	\hypN: X_a \sim P, X_b \sim P \quad
	\text{i.i.d. for some } P.
\end{align*}

\clearpage

\subsection{E Zoos}

\subsubsection{UI}

Based on the general case paper and the UI slide.

Let $\hat{\mu}_{\mid i-1}$ be any \emph{non-anticipating}\todo{What does it mean?}
estimator based on the first $i-1$ observations and we
denote the MLE estimator for the null as
$\hat{\mu}_{\mid n} = \argmax_{\mu\in\Theta_0} p_{\mu}(U^{{(n)}})$.

\begin{align*}
	S_{\textsc{UI}} := \frac{\prod_{i=1}^n q_{\hat{\mu}_{\mid i-1}} (U_{(i)})}{p_{\hat{\mu}_{\mid n}}(U^{(n)})}
\end{align*}

By i.i.d. assumption, the denominator reads the likelihood of $U^{(n)}$ given the MLE estimator.

It is mentioned that \emph{prequential ML} method is used here.

\Yongqi{Can I also used the MLE estimator in the numerator? I think so}

\subsubsection{RIPr}

By Theorem 1 in \cite{turnerExactAnytimevalidConfidence2023}, the GRO \E-variable
reads:
\begin{align*}
	S_{\textsc{RIPr}} := \frac{p_{\mu^*}(U^{(n)})}{p_{\mu^\circ}(U^{(n)})},
\end{align*}
where the prior put mass 1 on the boundary and $\mu^\circ = (n_a\mu^*_a/n + n_b \mu^*_b/n)$.

\begin{remark}
	In 2-stream Bernoulli and the null hypothesis is 2 group is distributed with the same distribution $P \in \Pc$,
	$S_{\textsc{cond}} = S_{\textsc{RIPr}} = S_{\textsc{PSEUDO}} = S_{\textsc{MIX}}$,
	which achieves GRO by corollary 1 of safe testing paper.
\end{remark}

\subsubsection{NML}

Also look at the luckiness

\subsubsection{COND}

Add a section on Wald's ratio test

We can first write the conditional \E-variable for a data block $X = U = (U_a, U_b)$:

\begin{align*}
	S_{\textsc{COND}} := \frac{q(X_a, X_b \mid X_a + X_b)}{p(X_a, X_b \mid X_a + X_b)}
	= \frac{q(X_a\mid X_a +X_b)}{p(X_a\mid X_a+X_b)}
	= \frac{f_{\textrm{FN}(n,n_a,Z,w)}(X_a)}{f_{\textrm{hyp}(n,n_a,Z)}(X_a)}.
\end{align*}

\begin{remark}
	Right now, I finally realize 4.2.3 of Yunda's thesis where the \E is:
	\begin{align*}
		S_{\textsc{COND}} := \frac{p_{\bm{\mu}}(X^{k-1}\mid Z)}{p_{\spri{\mu}}(X^{k-1}\mid Z)}.
	\end{align*}
	Here the $Z:=\sum_{j=1}^k X_j$ is the sum of the sufficient statistics and
	$X^{k-1}$ would be the sufficient statistics losing $1$ \emph{degree of freedom}.
\end{remark}

In our special case of Bernoulli streams, the conditional distribution corresponds to
the Fisher's noncentral hypergeometric (FN) distribution and the denominator is
simply the cases when odds ratio $\omega = 1$ (hypergeometric distribution).

\begin{align*}
	S_{\textsc{COND}} = \prod_{i=1}^n \frac{q(X_a^{(i)}\mid Z^{(i)})}{p(X_a^{(i)}\mid Z^{(i)})}
	= \prod_{i=1}^n \frac{f_{\textrm{FN}(n,n_a,Z^{(i)},w)}(X_a^{(i)})}{f_{\textrm{hyp}(n,n_a,Z^{(i)})}(X_a^{(i)})}
\end{align*}

\subsection{Calibrator}

First we run the standard Fisher's exact test and then obtain the \E-variable via
a calibrator $E = 1/\sqrt{q} - 1$.

\clearpage

\section{Questions and Comments}

\subsection{Prior or turning distribution}

We assume the oracle is some fixed point $\mu^*$
Let $\pi$ be the prior distribution on the alternative distributions.
The prior predictive distribution is defined as
\begin{align*}
	p(U) = \int p_{\mu}(U) \pi(\mu) \dd{\mu}.
\end{align*}

In the discussion with Ly on 02.05.2025, he pointed many names for
the above expression:
maximal likelihood distribution, the prior predictive distribution (Bayesian), and Bayes mixture (MDL).

Additionally, he mentioned that we usually used \textsc{letters} ($X,A$) for
observables where \textsc{Greek letters} are reserved for unknown
($\pi, \theta, \Theta$), especially in Bayesian framework.

In contrast to the belief's $\mu$ in Bayesian, MDL principle
regards $\pi$ on $\mu$ rather as a turning tool, uninformative,
to approach the oracle point $\mu^*$ given incoming data.

\subsection{Finding confidence sequence for case 1}
First of all, can we find CS for this testing setup.
If so, the confidence set should be searched through the $\Theta_0$
(dashed red line in \refFig{fig:case1_GRO}) or the entire $[0,1]^2$?

\subsection{Convexity of \textsc{SPACES}}

I am not sure why the modified settings where \cite{turnerExactAnytimevalidConfidence2023} and \cite{turnerGenericEvariablesExact2024} looks at the single outcome first and then extend to a single data block.

Calculating the KL divergence over the set of distributions $\mathcal{H}$ are convex while the calculation of KL over parameter space $\Theta$ is \textbf{ALSO} convex??

\begin{quote}
	% \textit{
	note that we cannot draw this conclusion if $\hypN'$ is not convex; for then the distribution $p_{W^\circ}'$ may not correspond to the distribution $p_{W^\circ}$ in the original problem — this correspondence is only guaranteed if $p_{W^\circ}'$ coincides with some $p_{\theta}'$.
	% }
\end{quote}

The part I am most confused about is the relationship between the convexity of the hypothesis space (original and modified) and parameter space. What we want is a $\hypN'$ that is compact and convex.

\href{https://statproofbook.github.io/P/kl-conv.html}{Proof: Convexity of the Kullback-Leibler divergence}, it gives the proof of convexity in the pair of probability distributions $(p, q)$ in discrete cases but the same could be applied to continuous cases with integration.

\subsection{What is effect size and why do we need it?}:
\cite{turnerGenericEvariablesExact2024} deals with effect size by stating that $\hypA$ consists of hypotheses with a minimal effect size $\delta$ while the parameter space of interest is $\hypN(\delta)$ (null hypothesis spaces) in \cite{turnerExactAnytimevalidConfidence2023}. These has confused me a great deal but I will try to explain in my words:

The effect size is the minimal difference between group a and group b
we would like to see when rejecting the null hypothesis.
Since we are concerned with a parametric model of Bernoulli stream in both group,
parameterized by $\theta_a$ and $\theta_a$.
The effect size would essentially be the divergence $d$ between $\theta_a$ and $\theta_b$
in the parameter space.

\subsection{Why puts a prior on the boundary only}

$W_1^*$: Prior on $\Theta(\delta)$; $W_1'$: Prior over $\Theta_1^>(\delta)$.
What are good reasons of putting all prior mass on the boundary?

Pros:
\begin{itemize}
	\item Computationally simpler as we only calculate once on $\theta_a$
	      and we get $\theta_b = \theta_a + \delta$ for free.
	\item If $(\theta_a^*, \theta_b^*)$ sits somewhere near the line,
	      \E-variable growth rate is higher with larger prior support.
	\item Theorem 2 in \cite{turnerExactAnytimevalidConfidence2023} guarantees a degenerate prior
	      on the boundary $\Theta(\delta)$
\end{itemize}

Cons:
\begin{itemize}
	\item If $d_{abs}(\theta_a^*, \theta_b^*) \gg \delta $,
	      \E-value calculated from $W_1'$ is much larger than $W_1^*$
\end{itemize}

\sout{What does it mean to \emph{pick $\hypA$ restricted to $\Theta(\delta)$}?}

In the discussion with Peter, he mentioned that the project $P^*$ might not be on the diagonal line unlike section \ref{sec:case1}.
It might just be a highly non-convex function that we can only solve by finding the min. of KL divergence after discretization.
He mentioned something about the projection not being straight but rather curve.

I think I made a mistake here, when looking at the divergent case,
the null would be the divergence bigger than $\delta$ while the alternative would be else?

{\color{red}
can we plug in $\theta_b^\circ = \theta_a^\circ + \delta$ here? If so,
it would just be solving a cubic equation of $\theta_a^\circ$.
}


\subsection{Re-orientate at 01.04.2025}

The biggest reason that I have been so disoriented is that I have so many documents, tests
scattered across platforms and places: \githubIssue, \personalRepo, \forkedRepo.
Originally, I wanna use Issue as a tracking board, personal repo for experiments,
and the forked repo for dev.
In the end, I became over-whammed by all and often found myself wasting energy,
organising thoughts even after just 1 day.

Thus, I have decided that I will put everything related to 2x2 in this \LaTeX\ document.
I will still use my \personalRepo for testing but I will not use \githubIssue any more.
I will write exclusively here as I find writing with \LaTeX\ is still the simplest
when even a little bit of math is involved.
I will use the item or numbered list here and \sout{strike through it} upon completion
or delete it when it is not longer relevant.

\subsection{Discussion with Dante at 23.04.2025}

The extended null hypothesis is meaningful now!

In this setup that I call ``one-shot'' where the whole data $U_a, U_b$ is seen at once,
the conditional \E-variable is optimal.

The extension of the null hypothesis is meaningful, in this parametric Bernoulli family only:
the extended null $\hypN'$ will be all distributions where all permutation $U$ are the same
and the marginal distributions of all (all exchangeable sequences).

In other words, let $U'=(U_1,U_2,\dots,U_{n_a+n_b})$ and
$\pi$ be any permutation operator on $\{1,\dots,n_a+n_b\}$.
The null hypothesis is now extended to a much bigger one,
namely, whether two streams $U_a, U_b$ are exchangeable,
\begin{align*}
	\hypN':=(U_1,U_2,\dots,U_{n_a+n_b}) \overset{d}{=}
	U_{\pi(1)},U_{\pi(2)},\dots,U_{\pi(n_a+n_b)}\}.
\end{align*}

Moreover, the marginal distribution of $U_j \overset{d}{=} U_i$ for any $i, j$
but the conditional dependence is not assumed in $\hypN'$ any more
in contrast to independence assumed in the original $\hypN$
where $U_i, U_j \iid P$.

Nevertheless, the $S_{\textsc{cond}}$ is still GRO-optimal, we can claim it is a
\textit{robust} \E-variable for a bigger problem.

\subsection{28.07.2025}

The goal to get a conditional \E-variable
\begin{itemize}
	\item with priors on log odds ratio: Jeffrey's, conjugate, Cauchy, Laplace, Gaussian
	\item uniformly most powerful (UMP) test for a fixed $\alpha$, ask Sebastian
	\item normalized maximal likelihood (NML), aligns with the MDL principle
\end{itemize}

The second amazingly achieves the most powerful test with a given test,
related materials are in the reply to the safe testing but it is given in
the context of unconditional one. So Peter was not sure about its look here.

Compare it with the old one (unconditional one)
\begin{itemize}
	\item Turner's
	\item Universal inference (MLE)
	\item p-to-e calibrator ($e = 1/\sqrt{q} - 1$ maybe)
\end{itemize}

Previous results shows that the UI performs the worst, next to Turner's,
a simple calibrator works best, e.g. $f(p) = \kappa p^{\kappa-1}$ for some $\kappa \in(0,1)$
and $f(p) = p^{-1/2}-1$.
noted that the Fisher's exact test gives a very small p-value

First step is to start with a fixed prior and fixed alternative.

\subsection*{Confusion about the sequential-ness and conditional-ness}

I was confused about the formulation between the sequential ratio and
the (what I called) global conditional \E-variables.

Let $X,Y$ be Bernoulli variables with $p$ and $q$, then we can write the
conditional likelihood ratio test (Wald's) as follow:

\begin{align*}
	R^n := \prod_{i=1}^n \frac{f_1(X_i, Y_i \mid X_i + Y_i)}{f_0(X_i, Y_i \mid X_i + Y_i)}
\end{align*}

and $f_1(\cdot) := (1 + \theta)^{-1}$ and $f_0 = (1 + \theta^{-1})^{-1}$

Here is the example of composite testing:
\begin{align*}
	\hypN: (X_1, Y_1), (X_2,Y_2),\dots \iid P \in \Pc_1 \: \textrm{v.s.}
	\hypA: (X_1, Y_1), (X_2,Y_2),\dots \iid P \in \Pc_{\theta}
\end{align*}

where $\Pc_1$ and $\Pc_{\theta}$ are distribution parametrized by odds ratio $\theta$.

Now it is almost the same as before but just extended where each data point $(X_1, Y_1)$
is replaced by a contingency table.
I think the reason I got so confused is here:
we are testing $m$ copies of tables $T_i$ and within that $T_i$,
each table is a 2x2 table of $n_a, n_b$ Bernoulli streams with $p, q$.

Using the notation $X^b_i, i = 1,\dots, n_a$ Ber and $X^a_i, i =1,\dots,n_b$
The full expression reads:
\begin{align*}
	\frac{f_1(X_1^a, X_2^a,\dots X_{n_a}^a, X_1^b, X_2^b,\dots X_{n_b}^b \mid Y_a + Y_b)}
	{f_0(X_1^a, X_2^a,\dots X_{n_a}^a, X_1^b, X_2^b,\dots X_{n_b}^b \mid Y_a + Y_b)}
\end{align*}

For a particular table $T_i$, I can just jump to the common doc, the hypothesis being tested is

\begin{quote}
	For all the 2x2 tables $T_i$, group a and group b frequency
	has no difference.
\end{quote}

\begin{align*}
	\hypN: T_1, T_2, \dots \iid P \in \Pc_1
	\quad \textrm{v.s.} \quad
	\hypA: T_1, T_2, \dots \iid P \in \Pc_{\theta}
\end{align*}

Questions to discuss tomorrow:

\begin{itemize}
	\item Is the \E-process-ness still here? \emph{YES}, by
	\item Is the optional continuation still here, I think so
	\item For the mixture, $\theta$ gets Bayesian updating each time (numerical)
	\item I think we assume group size $n_a, n_b$ is stationary right? Does the
	\item I remember there is some nice properties concerning the concave of the Fisher's information
\end{itemize}

\subsection*{Formulate this rigously: A small discussion with Peter about the general case paper}

If we have 1-d regular exponential family with the same sufficient statistics,
let's assume we are talking about Turner's example.

The GRO one is the conditional \E-variable and putting the prior on the curve (ark)
results also in the RIPr one (GRO) \E-variable.

Compared to this case (fixed odds ratio $\delta$ vs $1$, curve vs diagonal straight line),
it is not clear that if we still get the optimal \E-variable if instead of a fixed $\delta$,
we put a prior on the simplex, what sort of optimality we would get if there any.

Would the conditional variable still be the optimal one?


Let $f(k) = \binom{n_a}{k} \binom{n_b}{U-k}$ and $f'(k) = \frac{n_b!}{k!(n_a-k)!(U-k)!(n_b-U+k-n_a)!} = \binom{n_a - U + k}{n_a}$


\begin{align*}
	p_{n_a, n_b, u, \omega}(k) = \frac{\frac{n_b!}{k!(n_a-k)!(u-k)!(n_b-U+k-n_a)!}\omega^k}{\sum \frac{n_b!}{k!(n_a-k)!(u-k)!(n_b-U+k-n_a)!}\omega^j}
\end{align*}


We know $n_{11}$ is the $Y_a = k$ in our settings,





% \subsection{Walking through the maze of notation}

% Upon reading the literature, I find the notation between different papers quite different
% and it has been difficult for me personally to understand. 
% Here I will summarize the notations and their reasoning behind a particular representation.
% We denote calendar time index by $t$ to avoid the abuse of dots in sequential 
% or $k$-stream settings.

% In \cite{turnerExactAnytimevalidConfidence2023, turnerGenericEvariablesExact2024}, 
% sample follows some parameterized distributions 
% $P \in \Pc := \{P_{\theta}: \theta \in \Theta\}$. 
% The setup is restricted to 2 stream mostly where 
% $Y_{t,a} \iid P_{\theta_a}, Y_{t,b} \iid P_{\theta_b}$.
% % $\Theta_0, \Theta_1$ denote the parameter spaces of null and alternative hypothesis in $[0, 1]^2$. 

% In the $k$-stream sample test paper (\cite{haoEvalues$k$sampleTests2024} or 
% Ch. 4 in \cite{haoEvaluesAnytimevalidInference2025}). 
% It should be noted that the results are considered in a block, hence the E process-ness
% is OK in contrast to conditional and RIPr \E-variables in general paper.
% The null distributions $\Pc$ is some exponential family with mean-value parameter space $\Mc$.
% The null hypothesis \emph{relative to} $\Pc$ is 

% $\hypN$ is the collection of distributions (Bernoulli here)
% while $\Theta$ is parameter space in $[0, 1]^2$ 
% (whose convexity can simply visualized in the unit square).
% For example, $\Theta$ could be a convex set in the unit cube 
% while the distribution $\mathcal{H}$ or $\mathcal{P}$ might not be.

% In the case of Bernoulli distribution, which belongs to exponential family, 
% we denote the mean-value parameters space by $\Mc$ which contains all 
% $\mathtt{M} \in (0,1).$




% \textbf{Data}:
% A complete data block $Y$ of two streams $(Y_a, Y_b)$ is observed at each time-step, 
% indexed by the superscript $^{(t)}$: 
% \begin{align*}
%     &Y_a = (Y^{(1)}_a, Y^{(2)}_a, \dots, Y^{(m)}_a),\\
%     &Y_b = (Y^{(1)}_b, Y^{(2)}_b, \dots, Y^{(m)}_b).
% \end{align*}
% At time $t$, each data block $Y^{(t)} = (Y^{(t)}_{a}, Y^{(t)}_{b})$ 
% comprises two vectors of binary random variables of length $n_a$ and $n_b$
% , respectively, where all entries are i.i.d. according to the same Bernoulli distribution.
% \begin{align*}
%     &Y^{(t)}_{a,1}, \dots, Y^{(t)}_{a,n_a} \iid \mathrm{Ber}(\theta_a), \\
%     &Y^{(t)}_{b,1}, \dots, Y^{(t)}_{b,n_b} \iid \mathrm{Ber}(\theta_b).
% \end{align*}

% Within each group, $X^{(t)}_{g} := \sum_{i=1}^{n_g} Y^{(t)}_{g,i}, \; g \in \{a, b\}$ 
% is the sufficient statistic which is the number of ones in group g. 
% By default, We will restrict our discussion to the balanced design $n_a = n_b = 1$, 
% which greatly simplifies the setup into:
% \begin{align*}
%     Y^{(t)}_{a} \sim \mathrm{Ber}(\theta_a), \quad Y^{(t)}_{b} \sim \mathrm{Ber}(\theta_b).
% \end{align*}

% \clearpage

% \section{Case 1: Testing against an equal mean}
% \label{sec:case1}

% We are interested in whether two \textit{independent} Bernoulli streams 
% are equipped with equal means.
% We are testing if there is any discrepancy at all, hence a two-sided test.
% The GRO() \E-variable is the $S_{\textsc{mix}}$ below.

% $\hypN$, null hypothesis (composite):
% \begin{align*}
%     &\Theta_0 = \{ (\theta_a, \theta_b) \in [0, 1]^2 : \theta_a = \theta_b \} \\
%     &\hypN = \{P_{\theta_a, \theta_b} : (\theta_a, \theta_b) \in \Theta_0 \}
% \end{align*}

% $\hypA$, alternative hypothesis (simple):
% \begin{align*}
%     &\Theta_1 = \{ (\theta_a^*, \theta_b^*) \in [0, 1]^2 \} \\
%     &\hypA = \{P_{\theta_a, \theta_b}:  (\theta_a^*, \theta_b^*) \in \Theta_1 \}
% \end{align*}


% \begin{figure*}[ht!]
%     \centering
%     \begin{subfigure}[t]{0.49\linewidth}
%         \centering
%         \includestandalone{fig/Visual-unrestricted-projection}
%         \caption{The simple $\hypA$ is indicated by $C$ and 
%         the restricted parameter space of $\hypN$ is indicated by the 
%          dashed red line $AB$.
%         The RIPr of $\hypA$ onto $\hypN$ is denoted by the red dot ($C'$).}
%         \label{fig:case1_GRO}
%     \end{subfigure}%
%     ~ 
%     \begin{subfigure}[t]{0.49\linewidth}
%         \centering
%         \includestandalone{fig/Visual-restricted-abs}
%         \caption{Projection is skewed, as shown by the red dot.}
%     \end{subfigure}
%     \caption{Visualisation of a simple alternative against a composite null $\theta_a = \theta_b$}
%     \label{fig:GRO_projection}
% \end{figure*}

% \subsection{Euclidean distance and KL divergence in mean-value space} 

% Only in this special setup, the minimizer of Euclidean distance and KL divergence coincides.
% The $C'$ is the point on line $AB$
% that has the smallest Euclidean distance between $C$ ($CC' \perp AB$).
% \begin{align*}
%     C' = \inf_{P \in AB} d_2(P, C), \quad d_2(\cdot, \cdot) \textnormal{ is the Euclidean distance.}
% \end{align*}

% Luckily for us, it is exactly the same point RIPr.
% \begin{align*}
%     C' = \inf_{P \in AB} \KL{P}{C}, \quad \KL{\cdot}{\cdot} \textnormal{ is KL divergence.}
% \end{align*}

% In other composite settings, the RIPr of $\Qc$ onto $\Pc$ 
% \textit{does} falls on the boundary but unfortunately it is no longer the $l^2$ 
% projection 
% \footnote{
% $l^p$ usually refers to the sequence space where 
% $L^p$ are the function space. 
% }
% any more.

% \subsection{Case 1a: learn a simple alternative hypothesis}

% Does putting a prior turns this into a composite alternative hypothesis?

% \begin{remark}
%     We must not think of $\pi$ on $\Theta$ as a prior distribution but rather
%     a tuning parameter. Although the expression might resemble the Bayesian 
%     framework, the interpretation differs.

%     As suggested by Ly, one should not think of the simple alternative
%     $\Theta_1=\{(\theta_a^*,\theta_b^*)\in [0,1]^2\}$ with $\pi(\theta)$
%     as if we somehow \textbf{know} how $\pi$ distributes in $\Theta$.
%     Instead, $\pi$ is used here as a instrument to learn the true theta.
%     In simulation, we acted with known a fixed $(\theta_a,\theta_b)$.

%     The example Ly gave was very straightforward. In LASSO,
%     $\lambda$ is used to reduce the number of coefficients.
%     For example, the true model could be $y = \beta_3x_3 + \beta_{20}x_{20}$ with 
%     some Gaussian noise. Note that $\lambda$ is not contained in the true model
%     but only a medium for us to estimate the best model with corresponding constrains.

%     Therefore, we shall not say the simple alternative is extended to the composite setting
%     with a prior $\pi$. Instead, we are acting if we do not know the true alternative
%     and learn it via $\pi$.
% \end{remark}

% Theorem 2 of \cite{turnerExactAnytimevalidConfidence2023} stated that 
% if the point alternative is achieved by 
% $(\theta_a^*, \theta_b^*)$ if $(\theta_a^*, \theta_b^*) \in \Theta_0$, 
% meaning that the simple alternative sits within.

% Here we first write the KL between $(\theta_a, \theta_b) \in \Theta_0$ and a point alternative $(\theta_a^*, \theta_b^*)$:

% Eq. 10 in \cite{turnerExactAnytimevalidConfidence2023} is important as it states 
% that no matter which estimators we used, 
% we still get a \E-variable $\breve{\theta}_a \mid Y^{(j - 1)}$. 
% My understanding is that we will still get a \E-variable 
% as long as we condition on the previous filtration ($j - 1$) 
% but I am \textit{unsure} about the optimality about it.

% % \subsection{Case 1b: Extension to general null}

% % Modification: represented by a $'$ is single-outcome, changing from data block to a single event with $n_a / n$ of being group $a$ and $n_b / n$ of being group $b$.
% % \begin{equation*}
% %     \hypN' := \{P_{\theta}(G, Y) \mid \theta \in \Theta_0\}, \quad G = \{a, b\}
% % \end{equation*}

% % Construction is through a group variable $G$ but I am not sure why it would make for a convex $\hypN'$ in the case of a non-convex $\hypN$.

% \subsection{Finding the empirically best Beta prior}
% This is done in a minimax-sense: 
% one finds a $\gamma$ that is minimal among the worst case.
% We are basically looking for the smallest $\gamma$ in the worst case (minimax).
% But the problem is range of $\gamma$ could be vastly huge (unbounded) and 
% I have no idea how to locate one that is uniformly good across $\Theta$.
% % The worst case refers $\underset{(\theta_a, \theta_b) \in [0, 1]^2}{\max} E(\log()) - E()$.

% \subsection{Finding CS}
% The idea of CS is tight to the notion of effect size $\delta$,
% CS is simply calculated as case 2 with effect size $\delta = 0$.

% Let the confidence interval (CI) for level $1-\alpha$ at time $n$:
% \begin{align*}
%     \mathrm{CI}_{n, 1-\alpha} := \{\delta: S^{(n)}_{\delta} < 1 /  \alpha\},
% \end{align*}
% where $S_{\delta}^{(n)}$ is a \E-variable for null hypothesis $\Hc_{0,\delta}$,
% this ``good'' \E-variable is calculated via finding the minimizer of KL, discussed later.

% \begin{align*}
%     \Hc_{0, \delta} = \{P_{\theta_a, \theta_b}: d_{abs}(\theta_a, \theta_b) = \delta\}
% \end{align*}

% \clearpage

% \section{Case 2: Null hypothesis with linear boundary condition}

% Other names: absolute difference, absolute risk, linear boundary condition...

% For a divergence measure $d$ and a difference (effect size) $\delta$, 
% $\Theta_0(\delta)$ refers to the boundary condition.
% $\Theta_0^{<}(\delta)$ denotes the \textit{null} parameter spaces where $d \leq \delta$ 
% and $\Theta_1^{>}(\delta)$ denotes the \textit{alternative} parameter spaces where $d > \delta$.
% \Yongqi{
%     We ignore the boundary where $d = \delta$ as $P$ and $Q$ are absolutely continuous 
%     with respect to each other?
% }

% \textbf{Motivation}: Given a new drug, 
% instead of testing whether the control (a) and treatment (b) groups experience the same effect,
% a pharmaceutical company would like to determine 
% whether this drug performs at least $\delta$ better than the control group. 
% Here, $\delta$ represents the minimal relevant effect size.

% We are doing a one-sided test here to see if group b is \textit{at least} $\delta$ better than group a. 
% The ``better'' is measured in absolute difference in this case:
% \begin{align*}
%     d_{abs}(\theta_a, \theta_b) = \theta_b - \theta_a \in (-1, 1).   
% \end{align*}
% $d_{abs}$ in $[0, 1]^2$ ranges from $-1$ to $1$,
% we consider the case where $\delta > 0$ 
% as the same argument could be made for negative effect size by flipping the sign of $\delta$.
% We remove boundary cases $\delta = \{0, 1\}$ to improve numerical stability.

% $\hypN$, null hypothesis (composite): 
% \begin{align*}
%     &\Theta_0^<(\delta) = \{ (\theta_a, \theta_b) \in [0, 1]^2: d_{abs}(\theta_a, \theta_b) 
%     < \delta \}
% \end{align*}

% $\hypA$, alternative hypothesis restricted:
% \begin{align*}
%     &\Theta_1^>(\delta) = \{ (\theta_a, \theta_b) \in [0, 1]^2: d_{abs}(\theta_a, \theta_b) > \delta \}
% \end{align*}


% Theorem 2 in \cite{turnerExactAnytimevalidConfidence2023} tells us the RIPr \E-variable
% falls on the boundary domain.
% Fortunately, the \E-variable of highest power can be solved by finding the roots of
% the KL divergence, which is relatively easy here. 

% \subsection{Finding confidence sequence}

% Under $\alpha$ level and at time $m$, the exact (non-asymptotic) 
% confidence sequence $\CS{m}$ are the set of $\delta$ that 
% we have not been able to reject at time $m$ and level $\alpha$. 
% To obtained a two-sided CS, we simply give a list of,
% $\vec{\delta} = \{\delta_1, \delta_2, \dots, \delta_K \}$ ($K$: precision),
% and record at which $\delta'$ (effect size) the calculated cross $1/\alpha$.
% In the case of running intersection, we keep only the intersection.

% Differentiating KL against $\theta_a^\circ$:
% \Yongqi{This is very tricky as the differentiation depends on }
% \begin{align*}
%     \KL{\bm{\theta_1}}{\bm{\theta_0}}
%     &:=n_a\mathbf{E}_{Y\sim p_{\theta_a^*}}\left[\log\frac{p_{\theta_a^*}(Y)}{p_{\theta_a}(Y)}\right]+n_b\mathbf{E}_{Y\sim p_{\theta_b^*}}\left[\log\frac{p_{\theta_b^*}(Y)}{p_{\theta_b}(Y)}\right] \\
%     \dv{\mathrm{KL}}{\theta_a} 
%     &= \dv{}{\theta_a} 
%     n_a \qty(\theta_a^* \log \frac{\theta_a^*}{\theta_a} + (1-\theta_a^*) \log \frac{1-\theta_a^*}{1-\theta_a}) \\
%     & \quad + n_b \left(\theta_b^* \log \frac{\theta_b^*}{\theta_b} + (1-\theta_b^*) \log \frac{1-\theta_b^*}{1-\theta_b}\right) 
%      % \\
%      % &=  n_a \left( - \frac{\theta_a^*}{\theta_a} + \frac{1-\theta_a^*}{1-\theta_a}\right) 
%      % + n_b \cdot c \left(- \frac{\theta_b^*}{\theta_b} + \frac{1-\theta_b^*}{1-\theta_b}\right)
% \end{align*}

% \texttt{minimizeKL} gives us the solution of the following equation numerically,
% which is the derivative of the KL above\footnote{
%     If $\delta \in (-1, 0)$, calculation is over $(-\delta, 1)$ to ensure the parameters falls in $[0, 1]$
% }:
% \begin{align*}
%     n_{a} \left(-\frac{\theta_{a}^{*}} {\theta_{a}^{\circ}}+\frac{1-\theta_{a}^{*}} {1-\theta_{a}^{\circ}} \right) +
%     n_{b} \cdot c \cdot\left(-\frac{\theta_{b}^{*}} {\theta_{b}^{\circ}}+\frac{1-\theta_{b}^{*}} {1-\theta_{b}^{\circ}} \right)=0.
% \end{align*}



% \textit{Efficiency:} Solving $(\theta_a^\circ, \theta_b^\circ)$ becomes a cubic equation 
% when the solution only falls on the boundary by Theorem 2.
% It is 10x faster to use \verb|polyroot()| than \verb|stat::uniroot()| but 
% the trade-off is that the coefficients have to be provided for \verb|polyroot()|,
% which is bad for readability.

% Changes to be made:
% Use \verb|Re| to get the real roots of \verb|polyroot()| solution.



% \clearpage

% \section{Case 3: Null hypothesis with log odds ratio boundary condition}

% The distance here is calculated as 
% $d_{lOR}(\theta_a, \theta_b) = \log(\frac{\theta_b}{1 - \theta_b} \cdot \frac{1 - \theta_a}{\theta_a})$,
% the log odds ratio difference ranging from $-\infty$ to $\infty$.

% \subsection{Why CS is one-sided?}

% \begin{figure*}[h]
%     \centering
%     \begin{subfigure}[t]{0.49\linewidth}
%         \centering
%         \includestandalone[width=\linewidth]{fig/Ex-restricted-abs}
%         \caption{$d_{abs}(\theta_a, \theta_b) = \theta_b - \theta_a$}
%         \label{fig:abs diff}
%     \end{subfigure}%
%     ~ 
%     \begin{subfigure}[t]{0.49\linewidth}
%         \centering
%         \includestandalone[width=\linewidth]{fig/Ex-restricted-lOR}
%         \caption{$d_{lOR}(\theta_a, \theta_b) = \log(
%             \frac{\theta_b}{1 - \theta_b} \cdot \frac{1 - \theta_a}{\theta_a}
%             )$}
%     \end{subfigure}
%     \caption{
%     Examples of restricted alternative hypothesis parameter spaces for several values of two divergence measures.
%     $\Theta_0$ denotes the null hypothesis parameter space (dashed red line);
%     $\Theta_1^+(\delta)$, denoted by the shaded area,
%     is the restricted alternative hypothesis parameter space with difference larger than $\delta$ in $d_{abs}$ or $d_{lOR}$.
%     }
% \end{figure*}

% \clearpage

% \section{Case 4: Conditional \E-variable}
% \label{sec:cond}

% This is based on Dante's presentation on the Wald's RIPr (perturbed RIPr) and
% the Turner's RIPr along with the Example 9 in Ch. 4 in \cite{haoEvaluesAnytimevalidInference2025}.

% \textbf{Setup}: We have binary random variables $X, Y \in \{0, 1\}$, 
% $X \sim \mathrm{Ber}(\theta_a)$, $Y \sim \mathrm{Ber}(\theta_b)$.
% (We do not assume independent for now).
% \begin{align*}
%     \PR{X = 1} &= \theta_a, \, \PR{X = 0} = 1 - \theta_a \\
%     \PR{Y = 1} &= \theta_b, \, \PR{Y = 0} = 1 - \theta_b
% \end{align*}




% Let $u = \frac{\theta_a}{1-\theta_a}\big/\frac{\theta_b}{1-\theta_b} \in (0, \infty)$ 
% be the odds ratio between group a and b.
% It should be noted that $u$ along will not recover $(\theta_a, \theta_b)$ 
% hence we might need a nuisance parameter to recover CS.

% \begin{align*}
%     \condPR{X = 1, Y = 1}{X + Y = 2} &= 1 \\
%     \condPR{X = 1, Y = 0}{X + Y = 1} &= \frac{1}{1+u} \\
%     \condPR{X = 0, Y = 1}{X + Y = 1} &= \frac{u}{1+u} \\
%     \condPR{X = 0, Y = 0}{X + Y = 0} &= 1
% \end{align*}

% The key idea is \textit{enlarged} the distributions via reparameterization. 
% We can reparameterize the original distribution with $\alpha, \beta$
% to obtain the `marginalized' probability:
% \begin{align*}
%     \PR{X = 1, Y = 1} &= \alpha \\
%     \PR{X = 1, Y = 0} &= (1 - \alpha - \beta) \cdot \frac{1}{1+u} \\
%     \PR{X = 0, Y = 1} &= (1 - \alpha - \beta) \cdot \frac{u}{1+u} \\
%     \PR{X = 0, Y = 0} &= \beta
% \end{align*}

% \textit{Why the $\alpha, \beta$ and $\alpha', \beta'$ in two figures for perturbed RIPr?}
% This is to stress that the reparameterization can take any value.

% The original null hypothesis distribution family we consider is:
% \begin{align*}
%     \Pc = \bigcup_{u\in(0,\infty)} \Pc_{u}.
% \end{align*}

% Through reparameterization, we have the following enlarged family (\textit{proof needed}):
% \begin{align*}
%     \Pc_u^* = \{\mathrm{P}' \ll \Pc_u \mid 
%         \forall A \in \sigma: 
%         \condPR[\mathrm{P}']{A}{e} = \condERW{1_A}{e}, \mathrm{P}'\text{-a.s.}\}
% \end{align*}

% \Yongqi{What is the sigma algebra $\sigma$ and $e$ ?}

% The nice thing about this conditional is that we get a \E-variable that are close to 
% the GRO-optimal one but it works for any distribution constructed via the sufficient statistics
% while Turner's RIPr is only optimal in independent setting. 
% This is reminiscent of \cite{haoEvaluesAnytimevalidInference2025}'s argument 

% \emph{Why the mean of the NML is expected as 1?}

% The expectation is taken wrt the $Y_a$ and $Y_b$ follow the same Bernoulli? 
% Or is it composite with all same odds ratio?

% \begin{equation}
%     E_\textrm{NML}(Y_a) = \frac{1}{f_{\textrm{hyp}(n,n_a,u)}(Y_a)}\frac{f_{\textrm{FN}(n,n_a,u,\hat{w}(Y_a)}(Y_a)}{\sum_{k=\gamma_{\min}}^{\gamma_{\max}} f_{\textrm{FN}(n,n_a,u,\hat{w}(k))}(k)}
% \end{equation}

% Why this is \E-variable? Because the $E_{NML}$ would be just the NML density should 
% $Y_a$ follows the null hypothesis:
% \begin{align*}
%     \ERWi{Y_a \sim \hypN}{E_\textrm{NML}(Y_a)} 
%     &= \ERWi{Y_a \sim \hypN}{\condERWi{Y_a \mid U \sim \textrm{hyp}}{E_\textrm{NML}}{U}} \\
%     &= \ERWi{Y_a \sim \hypN}{\sum_{k = \gamma_{\min}}^{\gamma_{\max}} f_{\textrm{hyp}(n,n_a,u)}(k) 
%         \cdot E_\textrm{NML}} \\
%     &= \ERWi{Y_a \sim \hypN}{\sum_{k = \gamma_{\min}}^{\gamma_{\max}}
%         \frac{f_{\textrm{FN}(n,n_a,u,\hat{w}(k)}(k)}{\sum_{k=\gamma_{\min}}^{\gamma_{\max}} f_{\textrm{FN}(n,n_a,u,\hat{w}(k))}(k)} } = 1
% \end{align*}







% \subsection{Calculate stopping time and \E-value given confidence level and effect size}

% This is central to all the calculations and we call this function \verb|getStopE|.
% In safestats, it is called \verb|calculateSequential2x2E| with the arg \verb|simSettings == TRUE|.

% \begin{algorithm}
% \caption{Computation of \E-value given effect size and confidence level}
% \KwIn{Data $Y$; effect size $\delta$; confidence level $\alpha$; prior $W_a$ on $\theta_a$; grid size $G$}
% \KwOut{Stopping time $\tau$; final \E-value}

% Initialize $e_0, \tau, \E_{\tau}$\;

% Set a discrete uniform grid $\bm{g}$ on $(0, 1)$ of resolution $G$\;

% Set $\bm{g}_a \gets \bm{g} \cdot (1 - \delta)$\;

% Set $\bm{g}_b \gets \bm{g}_a + \delta$\;

% Compute $\theta_a \gets \bm{g}_a^\transpose \cdot W_{a}$\;

% Compute $\theta_a \gets \theta_a + \delta$\;

% Compute $\theta_0 \gets ({n_a} \cdot \theta_a + n_b \cdot \theta_b)/(n_a + n_b)$\;

% \For{$t \gets 1$ \KwTo $m$}{
%     Compute $e$ using likelihood ratio with $\theta$s\;

%     Update $e_t \gets e_{t-1} \cdot e$\;

%     \If{$e_t \geq 1 / \alpha$}{
%         Output $\tau \gets t$, $\E_\tau \gets e_t$\;
%         \textbf{break}\;
%     }

%     Update posterior density $W_a$ using $Y^{(t)}$\;

%     Compute $\theta_a \gets \bm{g}_a^\transpose \cdot W_{a}$\;

%     Compute $\theta_a \gets \theta_a + \delta$\;

%     Compute $\theta_0 \gets ({n_a} \cdot \theta_a + n_b \cdot \theta_b)/(n_a + n_b)$\;
% }

% \If(\tcp*[f]{No \E crossing $1/\alpha$}){$\tau = 0$}{
%     Output $\tau \gets m$, $\E_\tau \gets e_m$\;
%     \textbf{break}\;
% }

% \end{algorithm}



% It is generally the same with log odds ratio boundary condition except 
% \verb|stats::optim| is used to find $\theta_a^{RIPr}$.

% \begin{algorithm}
% \DontPrintSemicolon
% % set up aux. functions
% \SetKwFunction{calculateLikelihood}{calculateLikelihood}
% \SetKwFunction{updateBreveTheta}{updateBreveTheta}
% \SetKwFunction{minimizeKL}{minimizeKL}

% \caption{Calculation of confidence sequence for absolute difference constraint}
% \KwData{Data $Y$, confidence level $\alpha$, precision $K$, restricted effect size $\delta$,\\
%         \Indp prior $W_1^*$ on $\Theta(\delta)$: 
%         $W^*_{1,a} \sim \mathrm{Beta}(a_1, a_2), W^*_{1,b} \sim \mathrm{Beta}(b_1, b_2)$.}
% \KwResult{Confidence sequences, {\color{red}two-sided}}
% \BlankLine

% \For{$j = 1, 2, \dots, K$}{
%     % initialize $U_{0, a} = U_{0, b} \leftarrow 0$ (number of $1$s in each data block)\;
%     $e_0 \leftarrow 1$\;
%     $\breve{\theta}_a \leftarrow \updateBreveTheta(0, 0, a_1, a_2)$\;
%     $\breve{\theta}_b \leftarrow \updateBreveTheta(0, 0, b_1, b_2)$\;
%     $\theta_a^\circ \longleftarrow$ \minimizeKL{$\delta_j$, $\breve{\theta}_a$, $\breve{\theta}_b$}\;
%     $\theta_b^\circ \longleftarrow \theta_a^\circ + \delta_j$ \;
%     \For{$i = 1, 2, \dots, m$}{
%         $q_i \leftarrow 
%             \calculateLikelihood(Y^{(i)}_a, Y^{(i)}_b, \breve{\theta}_a, \breve{\theta}_b)$\;
%         $p_i \leftarrow 
%             \calculateLikelihood(Y^{(i)}_a, Y^{(i)}_b, \theta_a^\circ, \theta_b^\circ)$\;
%         $e_i \leftarrow e_{i-1} \cdot \frac{q_i}{p_i}$\;
%         % count $U_{i, a}, U_{i, b}$ in each block $Y_{i, a}, Y_{i, b}$\;
%         \If {$e_i \geq \frac{1}{\alpha}$ \textbf{and} running intersection}{
%             break, reject $\delta_j$\;
%         }
%         $\breve{\theta}_a \leftarrow \updateBreveTheta(
%             \sum_{j=1}^i Y_a^{(j)}, n_a\cdot i - \sum_{j=1}^i Y_a^{(j)}, a_1, a_2
%             )$\;
%         $\breve{\theta}_a \leftarrow \updateBreveTheta(
%             \sum_{j=1}^i Y_b^{(j)}, n_b\cdot i - \sum_{j=1}^i Y_b^{(j)}, b_1, b_2
%             )$\;
%         $\theta_a^\circ \longleftarrow$ \minimizeKL{$\delta_j$, $\breve{\theta}_a$, $\breve{\theta}_b$}\;
%         $\theta_b^\circ \longleftarrow \theta_a^\circ + \delta_j$ \;
%     }
% }
% \end{algorithm}


% \subsection{Calculating planned blocks with effect size and power}

% \begin{algorithm}
% \DontPrintSemicolon
% \SetKwFunction{minimizeKL}{minimizeKL}

% \caption{Calculation of nPlan given delta and beta}
% \KwIn{Data $Y$, effect size $\delta$, confidence level $\alpha$, prior $W_a$ for $\theta_a$, 
%     grid size $G$, \\ 
%     simulation time $M$, max stopping times $t_{max}$
% }
% \KwOut{currentWorstCaseQuantile}
% \BlankLine
% Set up a discretized grid for $\theta_a$ with $\delta, G$: $grid_a$\;
% Initialize worst case quantile: $q \leftarrow 0$ \;
% Initialize worst case power: $p \leftarrow 1$ \;
% \For{$t = 1, 2, \dots, G$}{
%     Initialize vectors of length $M$: stoppingTimes <- stopEs <- finalEs \;
%     $\theta_a \leftarrow grid_a[t]$ \;
%     $\theta_b \leftarrow fun(\theta_a)$ \;
%     \For{$i = 1, 2, \dots, M$}{
%         Sample $y_a$ of length $t_{max}$ from binomial distribution, $B(n_a, \theta_a)$\;
%         Sample $y_b$ of length $t_{max}$ from binomial distribution, $B(n_b, \theta_b)$\;
%         stoppingTimes <- stopEs <- finalEs <- getStopE(ya, yb)\;
%     }
%     Calculate current quantile: $q_t \leftarrow mean(stoppingTimes)$\;
%     Calculate current power $p_t$ as the percentages stopEs >= 1/alpha\;
%     \If{$p_t \leq p$}{
%         $p \leftarrow p_t$
%     }
%     \If{$q_t \geq q$}{
%         $q \leftarrow q_t$
%     }
% }
% \end{algorithm}
